\section{The local-global calculation}\label{sec:local-global}

The geometric work is now finished: the entire fiber is the root scheme of
one displayed polynomial.  The remaining argument is arithmetic, and it is
here that the factorization reveals why the family is so flexible.  The
cubic and quadratic do not both need a root at every place.  They only need
to cover all places between them.

Fix \(a\in\mathcal A\) and write
\[
 f_a(X)=(X^3-a)(X^2+X+1).
 \tag{3.1}
\]
The cubic has no rational root because \(a\) is not a cube, while the
quadratic is irreducible because its discriminant is \(-3\).  Hence
\(f_a\) has no rational root.  The fiber description
\eqref{eq:main-fiber} then gives
\[
 \Phi^{-1}(y_a)(\Q)=\varnothing.
 \tag{3.2}
\]

We now let the two factors divide the completions according to the residue
class of the prime.

\paragraph{The real place.}
The polynomial \(X^3-a\) has the real root \(\sqrt[3]{a}\).

\paragraph{The prime \(2\).}
Every prime divisor of \(a\) is \(1\pmod3\), so \(a\) is odd.  Thus \(X=1\)
is a root of \(X^3-a\) modulo \(2\), and its derivative \(3X^2\) is a
\(2\)-adic unit.  Hensel's lemma supplies a root in \(\Z_2\).

\paragraph{The prime \(3\).}
Write \(a=1+9k\) and substitute \(X=1+3T\).  Then
\[
 X^3-a=9h(T),\qquad
 h(T)=T+3T^2+3T^3-k.
 \tag{3.3}
\]
Modulo \(3\),
\[
 h(T)\equiv T-k,\qquad h'(T)\equiv1.
 \tag{3.4}
\]
The solution \(T\equiv k\pmod3\) lifts to a root in \(\Z_3\).

\paragraph{All other primes.}
Let \(p\ne2,3\).  If \(p\equiv1\pmod3\), then
\(\F_p^\times\) contains a nontrivial cube root of unity.  It is a simple
root of \(X^2+X+1\), so the quadratic has a root in \(\Q_p\).  This case
includes every prime dividing \(a\).

If \(p\equiv2\pmod3\), then \(p\nmid a\) and cubing is an automorphism of
\(\F_p^\times\).  The polynomial \(X^3-a\) therefore has a nonzero root
modulo \(p\).  Its derivative is nonzero, so the root lifts to \(\Q_p\).

This exhausts all places and proves the local assertions of
Theorem~\ref{thm:main}.  Notice that the same two factors alternate their
role with the residue class of \(p\).  The condition \(a\equiv1\pmod9\) is
needed only at \(3\); the prime-support condition controls the primes
dividing \(a\).  This is why primality disappears from the theorem: the
local proof remembers only a congruence and the support of \(a\).  That is
the arithmetic surprise behind the larger family.

\begin{remark}
The first prime parameter is \(a=19\), giving the classical
Berend--Bilu algebra \eqref{eq:berend-bilu}.  The first broader composite
example is
\(a=91=7\cdot13\).  The construction is therefore not merely a repackaging
of primes in one arithmetic progression.
\end{remark}

\begin{remark}
The reducibility of \(f_a\) is forced by the local condition.  A connected
finite \'etale \(\Q\)-scheme of degree greater than one has a transitive
Galois action containing a derangement.  Chebotarev then gives an unramified
prime at which the scheme has no point.  Hasse-failing zero-dimensional
schemes must consequently have at least two arithmetic components; see
Sonn \cite{Sonn} for the corresponding subgroup-covering formulation.
\end{remark}
